%% kapa: exam=1 round=26 year=2015 period=1 form=A subject=civil_law subject_ko=민법 seq=15 number=15
%% source: Q-Net 감정평가사 기출문제 (cst00309) artlSeq=5138001 file=extracted/1008/5138001/제26회_감정평가사 1차_기출/제26회_감정평가사 1차_기출/1교시_감평_A형.pdf
%% conversion: kordoc 4.24.51 -> tools/gen1.py (auto); review=proofread-sonnet; flags=none
\begin{question}{15}
무효와 취소에 관한 설명으로 옳지 않은 것은?
\begin{choices}
  \item 무효인 법률행위는 취소할 수 없다.
  \item 취소할 수 있는 법률행위의 추인은 취소의 원인이 소멸된 후에 하여야 효력이 있다.
  \item 「민법」상 법률행위의 일부가 무효인 때에는 전부를 무효로 함이 원칙이다.
  \item 취소할 수 있는 법률행위는 취소권자가 추인할 수 있고, 추인 후에는 취소할 수 없다.
  \item 취소권은 법률행위를 추인할 수 있는 날로부터 3년 내에 행사하여야 한다.
\end{choices}
\end{question}
\answer{1}
